\chapter{平台经济与双边市场}

\begin{examguide}
\textbf{本章考情}：双边市场理论是 0258 专业课的压轴考点（\important{专硕·核心}），Rochet–Tirole 定价公式的推导是各校计算题与论述题的命题素材；学硕微观命题中作为价格理论与垄断理论的新形式出现（\important{学硕·热点}）。本章是全书理论最深处，建议与第 2 章的古诺模型、第 3 章的交叉网络外部性对照学习。

\textbf{核心考点}：（1）双边市场的判定标准（价格结构非中性）；（2）交叉网络外部性与鸡生蛋问题；（3）Rochet–Tirole 定价公式的推导与倾斜式定价；（4）三种目标（利润、交易量、社会最优）下的定价比较；（5）Armstrong 竞争瓶颈模型与 Hotelling 双边扩展；（6）多归属与单归属决策；（7）负向网络外部性与平台规模的制约；（8）排他性协议的竞争效应与市场圈定的进入阻止条件；（9）平台包抄的机制与成立条件；（10）平台形态三阶段演进、平台生态圈与平台治理三内容（名解直考）。

\textbf{本章主线}：什么是双边市场（4.1）→ 平台如何定价（4.2）→ 平台如何竞争（4.3）→ 平台的策略行为（4.4）→ 平台如何演化为生态（4.5）。
\end{examguide}

\section{平台与双边市场概述}

\subsection{平台的经济学定义}

第 3 章末尾的论断构成了本章的起点：网络外部性使市场力量自我强化，这一力量既能成就卓越的平台，也能固化劣等的均衡。承载这一力量的组织形态就是平台，而平台面对的是两条相互缠绕的需求曲线：向买方收多少、向卖方收多少，两个价格必须同时决定，各自的效果又通过对方传导。这一结构使平台的定价、竞争与策略行为系统性地偏离单边市场的直觉，本章逐层拆解这套偏离，起点是界定平台本身。

\begin{definition}[平台（platform）]
平台是连接两个或多个相互依赖的用户群体、使它们之间发生直接互动的组织形态：平台自身不生产交易对象，只提供互动所需的基础设施与规则。
\end{definition}

与传统"管道式"企业（pipe）不同，管道式企业沿价值链单向传递价值（厂商生产、渠道分销、消费者购买），平台则退居幕后，让价值在用户之间直接产生：淘宝不生产商品，谷歌不生产内容，但二者都是各自市场中最有价值的组织。平台的这一性质使经济学分析的对象从"厂商—消费者"的单边关系转向"平台—两侧用户"的三方结构。

按连接对象与盈利来源，平台可分为四类。\textbf{交易平台}（transaction platform）撮合买卖双方完成交易，以佣金或交易服务费为主要收入，淘宝、美团、滴滴属此类；\textbf{广告平台}（advertising platform）以免费内容或服务聚集用户注意力，再把注意力售予广告主，谷歌、抖音属此类；\textbf{软件平台}（software platform）提供操作系统或开发框架，连接应用开发者与终端用户，iOS 与 Android 属此类；\textbf{支付平台}（payment platform）连接商户与消费者，以支付通道费等为盈利来源，支付宝与微信支付属此类。\mn{四类平台速记："交易促成交、广告卖注意、软件连开发、支付通两端"。四类的共同点是连接两侧、自身不生产交易对象，区别在于收费落在哪一侧，这正是 4.2 节倾斜式定价要解释的问题。}四类平台的差别集中落在收费结构上：交易平台向成交双方收费，广告平台对用户免费、向广告主收费，软件平台对用户低价、向开发者抽成，支付平台的收费方向随商户与消费者的议价能力而变。四类平台收费结构的差异服从同一套定价规律，4.2 节给出这套规律的完整推导。

\begin{figure}[htbp]
\centering
\begin{tikzpicture}[>=Stealth,
    plat/.style={rectangle, rounded corners=3mm, draw=main, ultra thick, fill=main!10, align=center, font=\small\bfseries, minimum width=22mm, minimum height=12mm},
    side/.style={rectangle, rounded corners=2mm, draw=second, thick, fill=second!8, align=center, font=\small, minimum width=30mm, minimum height=10mm},
    arr/.style={->, thick, second}]
% 节点坐标；连线放入 background layer 画在节点之下，节点最后绘制以盖住线头
\coordinate (P) at (0,0);
\node[side] (B) at (-3.8,1.9) {买方群体\\(消费者)};
\node[side] (S) at (3.8,1.9) {卖方群体\\(商家 / 司机 / 开发者)};
\begin{scope}[on background layer]
\draw[arr] (B) -- node[font=\small, color=structurecolor, above, sloped, pos=0.55] {$p_B$} (P);
\draw[arr] (S) -- node[font=\small, color=structurecolor, above, sloped, pos=0.55] {$p_S$} (P);
\draw[arr, main, ultra thick] (B) .. controls (-4.6,-2.0) and (4.6,-2.0) .. (S);
\draw[<->, thick, third, shorten <=-2pt, shorten >=-2pt] (B.east) -- (S.west); % 连接左框右边与右框左边，两端各插入框内 2pt，箭头主体露在框外
\end{scope}
\node[plat] at (P) {平台\\(淘宝 / 滴滴 / 微信)};
\node[font=\small\bfseries, color=main] at (0,-1.6) {交易互动（平台不直接参与）};
\node[font=\small, color=structurecolor, align=center] at (0,2.55) {交叉网络外部性：买方规模 $\uparrow$\\$\Rightarrow$ 卖方价值 $\uparrow$，反之亦然};
\end{tikzpicture}
\caption{双边市场的基本结构}
\label{fig:twosided}
\end{figure}

图~\ref{fig:twosided} 刻画了双边市场的基本结构。平台向两侧分别收费（$p_B$ 与 $p_S$），交易本身发生在用户之间；两侧的价值通过交叉网络外部性相互强化。这一结构决定了平台决策的基本形态：\important{平台必须同时决定两个价格}，且两个价格通过两侧需求的相互依赖而彼此关联——这正是本章定价分析的起点。

\subsection{双边市场的判定标准}

上一小节把平台界定为连接两侧用户的组织，但"连接两侧"这一描述性特征还不足以构成一个市场分类：百货商场同时面对供货商与顾客，传统媒体同时面对内容方与读者，它们是否都算双边市场？双边市场（two-sided market）理论的开创者 Rochet 与 Tirole 给出了一个可操作的判定标准：\important{价格结构非中性}（non-neutrality of the price structure）。这一标准把判断依据从"是否存在两侧用户"转到"价格在两侧之间如何分配是否影响交易量"，从而把双边市场从一般的多方交易关系中区分出来。

\begin{definition}[双边市场的判定标准（Rochet–Tirole）]
若市场中存在一个平台连接两类用户群体，且交易的\textbf{总量}不仅取决于平台向两侧收取的\textbf{价格总水平}（$p_B + p_S$），还取决于这一价格在两侧之间的\textbf{分配结构}（$p_B$ 与 $p_S$ 的相对大小），则该市场为\important{双边市场}；若交易量只取决于价格总水平而与分配结构无关，则该市场只是两个单边市场的简单加总。
\end{definition}

这一判定的学理含义在于：双边市场的本质特征是\important{两侧需求相互依赖}，平台对一侧降价的作用不只来自该侧的需求弹性，还通过交叉网络外部性传导到另一侧的需求。因此平台定价必须作为一个"结构"问题来处理：\important{总价格水平决定每笔交易的利润空间，价格分配结构决定交易规模}。举例而言，银行卡网络向持卡人与商户分别收取的费用如何分配，决定了银行卡的使用量：同样的总费用水平，向商户倾斜收费与向持卡人倾斜收费产生截然不同的交易规模，这正是价格结构非中性判定的现实含义。

中国银行卡的收费结构由监管直接规定，为价格结构非中性提供了一个制度层面的例证。按《关于完善银行卡刷卡手续费定价机制的通知》（国家发展改革委、中国人民银行，2016），借记卡的发卡行服务费费率不超过交易金额的 0.35\%、贷记卡不超过 0.45\%，网络服务费不超过 0.065\%，全部费用由商户一侧承担，持卡人刷卡不付费。监管在两侧之间选定的这一分配结构本身就是判定标准的应用：把费用集中于商户侧、令持卡人侧免费，才能维持持卡人的用卡规模，而用卡规模正是商户接受银行卡的理由。倘若把同样水平的总费用改由持卡人承担，用卡量会大幅萎缩，商户侧的价值随之下降，这正是"总水平不变、结构改变而交易量改变"的标准情形。

\subsection{交叉网络外部性与鸡生蛋问题}

价格结构之所以非中性，根源在于两侧需求被一种特定的外部性联系起来。第 3 章讨论的网络外部性发生在同一侧用户之间，双边市场的驱动力量则是它的跨侧形式：交叉网络外部性。这一概念是本章全部模型的共同基础，后续的定价公式、竞争均衡与策略分析都以它为出发点。

\begin{definition}[交叉网络外部性（cross-side network externality）]
交叉网络外部性是指平台一侧用户规模的变化直接改变另一侧用户效用的外部性：一侧用户规模的扩大提升另一侧用户的价值，两侧用户的效用函数中相互包含对方的规模。
\end{definition}

买家越多，商家入驻的价值越高；商家越多，买家逛平台的价值越高。\mn{交叉网络外部性与直接网络外部性（第 3 章）的区别：直接外部性发生在\textbf{同一侧}用户之间（微信好友），交叉外部性发生在\textbf{两侧}用户之间（商家与买家）。}两侧互为对方的"产品价值"来源，这一相互依赖产生了平台启动的经典难题，即\textbf{鸡生蛋问题}。

\begin{definition}[鸡生蛋问题（chicken-and-egg problem）]
鸡生蛋问题是指双边市场启动阶段的互为因果困境：买家不来因为商家少，商家不来因为买家少，两侧都不愿率先进入，市场无法从零规模自发启动。
\end{definition}

鸡生蛋问题的严重程度可以精确刻画。把两侧的参与决策写成依赖对方规模的条件，就能看清市场为什么会停在零规模，以及需要多大的初始规模才能脱离这一状态。下面的推导给出零均衡的自我实现机制与两侧的临界规模门槛。

\begin{derivation}{鸡生蛋问题的形式化}{}
设商家入驻的价值为 $v_S(n_B)$，买家接入的价值为 $v_B(n_S)$，二者均为递增函数，且 $v_S(0) = v_B(0) = 0$：没有另一侧用户时，平台对任何人都没有价值。两侧的参与条件分别为
\[
v_S(n_B^e) - p_S \geq 0, \qquad v_B(n_S^e) - p_B \geq 0,
\]
其中上标 $e$ 表示\textbf{预期}：用户根据对另一侧规模的预期做决策，均衡要求预期自我实现（第 3 章）。

第一步，检验零规模能否自我维持。在启动阶段 $n_B^e = n_S^e = 0$，由 $v_S(0) = v_B(0) = 0$，两侧的参与条件在 $p_S > 0$、$p_B > 0$ 时均不成立，于是两侧都不参与，实际规模 $n_B = n_S = 0$ 恰与预期吻合。\important{$(0,0)$ 因此是一个自我实现的均衡}，市场无法从零规模自发启动。

第二步，求正规模均衡所需的条件。由 $v_S$ 递增，参与条件 $v_S(n_B^e) \geq p_S$ 等价于 $n_B^e \geq v_S^{-1}(p_S)$，买方一侧同理。记
\[
\underline{n}_B = v_S^{-1}(p_S), \qquad \underline{n}_S = v_B^{-1}(p_B),
\]
二者分别是使卖方愿意入驻的最小买方规模与使买方愿意接入的最小卖方规模。任何正规模均衡 $(n_B^*, n_S^*)$ 必须同时满足
\[
n_B^* \geq \underline{n}_B, \qquad n_S^* \geq \underline{n}_S.
\]
这一对不等式界定了双边市场的\important{临界规模}：两侧规模必须同时越过各自的门槛，任一侧跌破门槛都会使另一侧退出，系统退回零均衡。

第三步，考察价格如何进入门槛。$v_S^{-1}$ 与 $v_B^{-1}$ 同为递增函数，故 $\partial \underline{n}_B/\partial p_S > 0$、$\partial \underline{n}_S/\partial p_B > 0$，即提价抬高对方一侧的门槛、降价压低门槛。特别地，当 $p_S \leq 0$ 时，由 $v_S(0) = 0 \geq p_S$ 可知 $\underline{n}_B = 0$，\important{负价格把该侧门槛完全消除}：卖方在买方规模为零时也愿意入驻。
\end{derivation}

上述条件把平台启动问题转化成一项明确的任务：或者把两侧规模同时推过门槛，或者把门槛本身降到零。现实中的破解方式对应三条路径。其一，\textbf{补贴一侧}：平台对一侧定价为负，发补贴、免佣金，按第三步的结论令该侧门槛归零，以人工制造的初始规模撬动另一侧进入，正反馈随之启动，外卖平台补贴用户、支付平台补贴商户都属此类。其二，\textbf{平台自营一侧}：平台自己充当一侧用户，以自营商品或自有内容直接供给初始规模，京东以自营起家、视频平台以自制内容开局都循此路。其三，\textbf{渐进式接入}：先以单边工具积累一侧用户，再开放另一侧，微信先做通讯工具、再开放公众号与小程序即为典型。三种策略的共同逻辑是人为创造初始规模的"种子"，越过第 3 章所述的临界规模；它们付出的代价则各不相同：补贴付出的是现金，自营付出的是资产与库存周转，渐进式接入付出的是时间。

\textbf{案例：淘宝以免费策略击溃 eBay 易趣（2003–2006）。}淘宝与 eBay 易趣的竞争是中国平台经济的第一场标志性战役，其胜负机制正是上述门槛条件的现实演示。2003 年淘宝上线时，易趣已占据中国 C2C 市场的绝大部分份额，且按 eBay 全球模式向卖家收取上架费与交易佣金。淘宝的策略是双重的：其一，\textbf{对卖家侧完全免费}，卖家零成本开店，对应上式中 $p_S \leq 0$ 的情形，卖方侧门槛 $\underline{n}_B$ 被直接压到零，易趣多年积累的商家池壁垒随之失效；其二，\textbf{配套即时通讯工具（阿里旺旺）降低买卖双方的互动成本}，抬高 $v_S$ 与 $v_B$ 两条价值曲线，使同样的规模能够支撑更强的参与意愿。按 4.2.3 节倾斜式定价理论，淘宝的免费策略是把价格结构向卖家侧大幅倾斜：卖家侧的需求弹性高（新平台对卖家缺乏锁定）、卖家侧对买家侧的交叉外部性贡献大（丰富的商品供给是吸引买家的前提），补贴卖家以撬动买家侧规模。易趣受制于全球统一的收费模式无法跟进，商家与买家同步迁移，2006 年易趣中国业务实质上退出市场。这一案例的完整理论含义在于：\important{平台竞争的决定性变量是价格结构而非价格总水平}。淘宝与易趣的"总收费"差异远小于"结构"差异，而结构差异足以决定平台的生死。

\section{双边市场基准模型：Rochet–Tirole 模型}

\subsection{模型设定}

上一节的门槛条件说明了平台为什么要补贴，却没有回答补贴给哪一侧、补到什么程度。回答这两个问题需要一个完整的定价模型。Rochet 与 Tirole 的基准模型抓住了双边市场定价的本质：平台对两侧分别收取交易费，交易量由两侧需求共同决定，平台在两个价格之间权衡。这是本章的核心推导，也是各校计算题最常见的命题来源。

\begin{derivation}{双边市场定价公式}{}
设平台连接买方（$B$）与卖方（$S$），平台对每笔交易向买方收取 $p_B$、向卖方收取 $p_S$，平台提供每笔交易的边际成本为 $c$。设两侧的交易参与率分别为 $D_B(p_B)$ 与 $D_S(p_S)$（即价格上升时该侧参与比例下降），且两侧的参与共同决定交易量：
\[
V = D_B(p_B) \cdot D_S(p_S).
\]
这一交易量函数的含义是：交易需要买卖双方\textbf{同时}参与，两侧参与率的乘积刻画"双方都在场"的概率或比例。

平台的利润为
\[
\pi = \bigl(p_B + p_S - c\bigr) \cdot D_B(p_B) \cdot D_S(p_S).
\]
对 $p_B$、$p_S$ 求一阶条件：
\[
\frac{\partial \pi}{\partial p_B} = D_BD_S + (p_B+p_S-c)\,D_B'D_S = 0
\;\Longrightarrow\; p_B+p_S-c = -\frac{D_B}{D_B'},
\]
\[
\frac{\partial \pi}{\partial p_S} = D_BD_S + (p_B+p_S-c)\,D_BD_S' = 0
\;\Longrightarrow\; p_B+p_S-c = -\frac{D_S}{D_S'}.
\]
记两侧需求弹性 $\eta_i = -D_i'\,p_i/D_i > 0$（价格上升 $1\%$ 时该侧参与下降的百分比），由 $-D_i/D_i' = p_i/\eta_i$，上式化为
\[
p_B+p_S-c = \frac{p_B}{\eta_B}, \qquad p_B+p_S-c = \frac{p_S}{\eta_S}.
\]
两式相等即得利润最大化的\textbf{价格结构条件}
\[
\frac{p_B}{\eta_B} = \frac{p_S}{\eta_S}.
\]
由合比性质，$\frac{p_B}{\eta_B} = \frac{p_S}{\eta_S} = \frac{p_B+p_S}{\eta_B+\eta_S}$。记总价格水平 $M = p_B+p_S$，代入价格总水平等式：
\[
M - c = \frac{M}{\eta_B+\eta_S}
\;\Longrightarrow\; \frac{M-c}{M} = \frac{1}{\eta_B+\eta_S},
\]
解得
\[
M = \frac{\eta_B+\eta_S}{\eta_B+\eta_S-1}\,c,
\]
其中最后一步要求 $\eta_B+\eta_S>1$（两侧弹性之和大于 $1$ 时内点解存在，与单边垄断要求 $\eta>1$ 同构）。再按弹性占比分摊，得到各侧价格：
\[
p_i = \frac{\eta_i}{\eta_B+\eta_S}\,M = \frac{\eta_i}{\eta_B+\eta_S-1}\,c, \qquad i = B,S.
\]
\end{derivation}

两个等式分别揭示了双边市场定价的两个核心结论。\mn{Rochet–Tirole 定价公式的记忆线索：单边市场勒纳公式是 $p-c = p/\eta$；双边市场由 $M-c = M/(\eta_B+\eta_S)$ 确定总价格水平，再按 $p_B/\eta_B = p_S/\eta_S$ 分摊到两侧。一句话：弹性占比决定分摊，弹性大的一侧承担总价格中的更大份额。}\textbf{价格结构规则}：$\frac{p_B}{\eta_B} = \frac{p_S}{\eta_S}$ 表明，利润最大化的价格结构使两侧的\important{价格-弹性比相等}。这一条件的直观含义是：总价格水平给定时，把 1 元价格从一侧挪到另一侧，两侧参与率的百分比变动分别为 $\eta_B/p_B$ 与 $\eta_S/p_S$；两者相等意味着任何再分摊都不改变交易量，平台处于"再调整价格结构已无利可图"的临界点。封闭解 $p_i = \frac{\eta_i}{\eta_B+\eta_S-1}c$ 给出更直接的读法：\important{弹性较大的一侧承担总价格中更大的份额}。这是本模型（对交易收费、两侧需求仅由本侧价格决定）的特征结论，与"弹性大的一侧获补贴"的经验印象不同，后者的模型基础是交叉外部性显式进入效用的 Armstrong 框架（4.3.2 节）。\textbf{价格总水平规则}：总价格水平 $M$ 满足 $\frac{M-c}{M} = \frac{1}{\eta_B+\eta_S}$，价格加成由两侧弹性之和决定，这与单边市场勒纳公式（Lerner formula）$p = c \cdot \eta/(\eta-1)$ 的结构一致：单边垄断的加成由单一市场的弹性决定，双边平台的加成由两侧弹性之和决定；各侧价格按弹性占比分摊总水平，分摊结构由两侧弹性之比决定。

把上述模型结果翻译成经济学语言，双边市场定价有三个层次的结论。其一，\important{定价决策的本质是选择结构而非选择水平}：两个平台的总收费水平相同，利润与交易量却可以截然不同，差别全在两侧价格的分配上，这正是淘宝击败易趣时发生的事。其二，\important{补贴是定价的自然结果而非慈善}：基准封闭解中两侧价格均为正，低于边际成本甚至为负的价格来自两处。一处是本节的便利收益变体：卖方承担的净费用一旦超过净成本，买方侧价格即转负（$p_S > c-b$ 时 $p_B < 0$），买方侧的负价格实为卖方便利收益的返还；另一处是交叉外部性显式进入效用的 Armstrong 框架：被另一侧强烈需要的一侧获得净补贴（4.3.2 节）。平台在低价侧的每一分亏损都是对另一侧规模的投入。其三，\important{免费模式的根源在价格结构倾斜}：数字平台普遍采取"一侧免费、另一侧收费"的结构，免费侧凭借其对另一侧的外部性贡献充当撬动另一侧的杠杆（4.3.2 节），免费本身正是这一杠杆作用的定价表现。\mn{三个结论速记：结构重于水平；补贴是投资；免费源于倾斜结构。}这三个结论把定价公式与银行卡网络交换费、电商平台补贴买家抽佣商家、操作系统对用户免费等经验现象直接对应起来，具体展开见 4.2.3 节。

基准模型假定两侧从交易中获得的收益已完全反映在各自的需求函数中。银行卡网络一类市场需要一个直接的推广：商户接受银行卡本身能带来销量提升与现金管理成本的节省，这部分收益并不经由平台收费体现。设卖方每笔交易额外获得\textbf{便利收益}（convenience benefit）$b$，则卖方的净支付为 $p_S - b$，卖方需求相应写作 $D_S(p_S - b)$。以 $\tilde p_S = p_S - b$ 记卖方净价格，平台利润为
\[
\pi = \bigl(p_B + \tilde p_S + b - c\bigr)D_B(p_B)D_S(\tilde p_S) = \bigl(p_B + \tilde p_S - (c - b)\bigr)D_B(p_B)D_S(\tilde p_S),
\]
这与基准模型的形式完全一致，只是边际成本 $c$ 被\important{净成本 $c-b$} 取代。因此前述全部结论原样成立，只需处处以 $c-b$ 替换 $c$。把买方侧的一阶条件 $p_B + p_S - (c-b) = p_B/\eta_B$ 整理成 $p_B$ 的显式表达，得到考试中最常直接考查的形式
\[
p_B = \frac{\eta_B}{\eta_B - 1}\,(c - b - p_S).
\]
\mn{便利收益 $b$ 的记法：$b$ 是卖方"白得"的好处，因而抵消平台的边际成本，公式里 $c$ 一律换成 $c-b$。由此可一眼看出补贴条件：$p_S > c-b$ 时买方侧价格为负。}这一表达给出了买方侧补贴的判定条件：$\eta_B > 1$ 时系数为正，故\important{当 $p_S > c-b$ 时 $p_B < 0$}，即卖方侧承担的费用一旦超过净成本，平台就应当对买方侧发放净补贴；便利收益 $b$ 越大，卖方越愿意承担费用，买方侧获得补贴的空间也越大。

银行卡网络是这一变体的标准应用。商户接受银行卡获得的便利收益包括客单价提升、收银效率改善与现金保管成本下降，数额可观。按上式，$b$ 较大使净成本 $c-b$ 很低甚至为负，商户侧费用 $p_S$ 轻易超过净成本，持卡人侧因此获得负价格。这正是 4.1.2 节所述中国刷卡手续费结构的理论解释：全部费率由商户承担，持卡人不但免费，还常常通过信用卡积分与返现获得实质上的负价格。国际卡组织的\textbf{交换费}（interchange fee）制度同理，由收单机构代商户向发卡机构支付，用于补贴持卡人一侧。交换费水平因此成为各国监管的焦点：费率过高会抬高商户成本并最终转嫁给全体消费者，过低则削弱发卡机构补贴持卡人的能力、抑制用卡规模。这一权衡是双边市场监管中最经典的难题，也是欧盟与美国先后立法限制交换费的由来。

\subsection{三种目标的定价比较}

平台的定价目标不限于利润最大化。比较三种目标下的定价结构，可以揭示平台行为偏离社会最优的方向。比较需要一个统一的福利度量，此处给出全书通用的福利分析基本方法。\textbf{消费者剩余}（consumer surplus）度量消费者从交易中获得的总收益：对每一单位交易，支付意愿减去实付价格即该单位的剩余，全部成交单位的剩余总和就是反需求曲线与价格线之间的面积。设需求函数为 $D(p)$，价格为 $p$ 时的消费者剩余为
\[
CS(p) = \int_{p}^{\bar p} D(s)\,ds,
\]
其中 $\bar p$ 为使需求降为零的价格上界。由微积分基本定理，$dCS/dp = -D(p)$：价格每上升一单位，消费者剩余减少当前需求量（全部存量消费者各多付一元，加上边际消费者退出）。\textbf{生产者剩余}（producer surplus）是生产者从交易中获得的总收益，在平台定价问题中即平台利润。\textbf{社会总福利}（total surplus）是消费者剩余与生产者剩余之和，对双边平台即两侧用户剩余加平台利润 $W = CS_B + CS_S + \pi$；价格支付是消费者与平台之间的转移，相互抵消，不直接进入总福利。\textbf{社会最优}指最大化 $W$ 的价格结构，它是评价利润最大化定价与交易量最大化定价的基准。

\begin{derivation}{三种定价目标}{}
继续使用 $V = D_BD_S$ 的设定。\textbf{交易量最大化}的平台最大化 $V$。由于 $D_B' < 0$、$D_S' < 0$，交易量对两个价格都严格递减，内部解不存在：交易量最大化的平台会持续降价直至两侧参与饱和（$D_B = D_S = 1$）。以线性参与函数 $D_i = 1 - a_ip_i$（参与率上限为 $1$）为例，饱和要求 $p_i \le 0$，故交易量最大化的定价不高于零（免费甚至补贴），低于利润最大化定价。

\textbf{社会最优}的平台最大化总剩余：两侧用户剩余加平台利润，
\[
W = CS_B(p_B) + CS_S(p_S) + (p_B+p_S-c)\,D_BD_S,
\]
其中 $CS_i(p_i) = \int_{p_i}^{\bar p_i} D_i(s)\,ds$ 为侧 $i$ 的用户剩余（反需求曲线下的面积，上文已给出定义），$\bar p_i$ 为使参与降为零的价格上界。利用 $\frac{dCS_i}{dp_i} = -D_i(p_i)$，一阶条件为
\[
\frac{\partial W}{\partial p_B} = -D_B + D_BD_S + D_B'D_S(p_B+p_S-c) = 0,
\]
解得
\[
p_B+p_S-c = \frac{D_B(1-D_S)}{D_B'D_S} < 0,
\]
其中不等号来自 $D_B' < 0$ 且参与不完全时 $D_S < 1$（对 $p_S$ 的一阶条件对称）。社会最优的\textbf{总价格水平低于边际成本}：交易需要两侧同时在场，价格高于边际成本时，一侧用户的接入价值因另一侧缺位而无法实现的部分（分子 $D_B(1-D_S)$）构成纯损失，社会最优以低于成本的定价（补贴交易）压缩这一损失，直至补贴扩大的交易与边际无谓损失相抵；补贴的强度随另一侧的参与缺口（$1-D_S$、$1-D_B$）扩大而提高。当两侧参与饱和（$D_B = D_S = 1$）时，条件退化为标准的边际成本定价 $p_B+p_S = c$。

三种目标的价格排序为
\[
p^{\text{利润最大化}} > p^{\text{社会最优}} > p^{\text{交易量最大化}}.
\]
以对称线性参与函数 $D_i = 1-p_i$ 为例可以验算：利润最大化由 $M-c = M/(\eta_B+\eta_S)$ 给出，代入 $M = 2p$、$\eta_i = p_i/(1-p_i)$ 得 $2p-c = 1-p$，即各侧价格 $p = (1+c)/3$；社会最优由 $2p-c = -p$ 给出 $p = c/3$；交易量最大化价格为 $0$。三者满足 $(1+c)/3 > c/3 > 0$，与上述排序一致。
\end{derivation}

这一排序的政策含义是直接的。竞争性的平台被市场压力推向交易量最大化的方向，其定价低于垄断平台的利润最大化定价，因此\important{平台竞争本身就在保护消费者}；而利润最大化的垄断平台相对社会最优\important{过度收费}，其价格结构同时偏离了社会最优所要求的补贴结构。两重偏离叠加，构成平台监管干预的福利理论基础，第 6 章的反垄断分析正建立在这一基础之上。这里还有一个最常被考查的辨析点：社会最优要求总价格水平低于边际成本，即对平台交易给予\important{净补贴}，这与单边市场"价格等于边际成本"的教科书结论明显不同，根源在于双边市场的交易需要两侧同时在场，一侧的参与不足会使另一侧的接入价值落空，补贴正是用来压缩这部分落空的损失。

\subsection{倾斜式定价}

利润最大化的价格结构规则直接导出双边市场最显著的经验特征，即倾斜式定价。前两小节的公式已经蕴含了这一特征，此处单独提出有两重理由：它是双边市场在现实中最容易观察到的现象，又极易与掠夺性定价混淆，因而成为辨析题的常客。

\begin{definition}[倾斜式定价（skewed pricing）]
倾斜式定价是指平台对一侧用户收费极低甚至为负（补贴）、对另一侧收费较高的定价结构：价格结构向弹性大、外部性贡献大的一侧大幅倾斜。
\end{definition}

补贴哪一侧的规则由两个因素决定。其一，该侧对另一侧的\textbf{外部性贡献}：该侧参与给另一侧带来的收益越大，平台越应补贴该侧以撬动另一侧的规模。这一机制在 4.2 节的便利收益变体与 4.3.2 节的 Armstrong 模型中分别有明确的数理形式：买方参与给卖方带来便利收益 $b$，$b$ 进入净成本 $c-b$ 压低了买方侧价格（$p_S > c-b$ 时买方侧价格转负）；卖方强烈需要买方时 $\alpha_S$ 大，Armstrong 模型给出买方侧低价。两个框架方向一致：\important{给另一侧带来收益多的一侧获得补贴}。其二，该侧的\textbf{需求弹性}：经验上被补贴的一侧通常也是弹性大的一侧，弹性大意味着降价撬动的参与多，补贴的杠杆效果强；但 4.2 节基准封闭解显示，当两侧需求互不进入对方效用时，弹性大的一侧反而承担总价格中更大的份额，因此"弹性大的一侧获补贴"应理解为经验规律与 Armstrong 框架的结论，而非基准模型的推导结果。考题通行的标准表述是"弹性大、外部性贡献大的一侧获得补贴"，作答时以这一口径为准。外卖平台的实践完全符合这一逻辑：消费者侧的流量是商家的生命线（外部性贡献大）、对价格敏感（经验上弹性大），故补贴消费者；商家侧依赖平台流量，故承担佣金。\mn{倾斜式定价案例：外卖平台补贴消费者、向商家抽佣；游戏主机低价出售主机（亏损）、向游戏抽成；支付平台对商户免费、向消费者收费或反之。判断"补贴哪一侧"的规则：对另一侧外部性贡献大的一侧获得补贴（经验上该侧通常也弹性大），考试标准表述："弹性大、外部性贡献大的一侧获得补贴"。}价格结构"一侧免费、另一侧付费"是利润最大化定价的自然结果，与平台是否"慷慨"无关：免费侧的亏损构成撬动付费侧的投入成本。

\begin{misconception}
\textbf{误区}：把倾斜式定价等同于掠夺性定价（predatory pricing）。二者表面相似（低价甚至负价），性质不同：倾斜式定价是双边市场\textbf{结构}的均衡结果，低价侧与另一侧的高价同时存在、相互配套，长期持续；掠夺性定价是单边市场的策略性行为，为排挤对手而低于成本销售，待对手退出后提价回收。判定要点：其一，低价是否与另一侧收费构成交叉补贴结构（倾斜式定价的特征）；其二，低价是否为持续性均衡行为而非阶段性排挤。二者并非互斥，补贴战同时具有排挤动机，双边市场下掠夺性定价的判定（成本与回收如何界定）是第 6 章反垄断分析的难题。
\end{misconception}

\textbf{案例：滴滴与快的的补贴大战（2014）。}2014 年，滴滴与快的在网约车市场展开的补贴战是中国互联网史上规模最大的倾斜式定价竞赛。两家平台同时对乘客与司机双向补贴：乘客每单立减、司机接单奖励。以双边市场框架解读：司机侧与乘客侧的交叉外部性极强（车多乘客就愿意等，乘客多司机就有单接），而两侧的需求弹性都高（用户多归属、转换成本低），按上述补贴规则，两侧都应当获得补贴，平台利润只能来自未来赢家通吃后的佣金回收。双方在数月内投入的补贴规模以十亿元计。补贴战的终局是 2015 年 2 月两家宣布合并，市场从双寡头补贴竞争转入垄断定价阶段：补贴逐步退坡、抽成逐步上调。\important{补贴、合并、提价的三部曲此后在中国平台经济反复出现}，共享单车与社区团购皆循此路，其机制正是本章理论所刻画的：网络外部性市场的前期投入以竞争性补贴的形式支出，垄断格局形成后再以佣金的形式回收。

\section{双边市场的竞争模型}

\subsection{多归属与单归属}

前两节把平台当作单一垄断者处理，价格结构的全部约束来自两侧的需求弹性。市场上一旦出现第二个平台，约束就改变了：用户可以在平台之间择一，也可以同时使用多个平台。平台竞争分析的起点正是用户的这一选择，它决定了竞争的激烈程度，也决定了价格结构倒向哪一侧。

\begin{definition}[单归属与多归属]
\textbf{单归属}（single-homing）指用户只加入一个平台（如多数用户的微信）；\textbf{多归属}（multi-homing）指用户同时使用多个平台（商家同时入驻淘宝与京东、消费者手机上同时装有多个外卖 App）。
\end{definition}

用户选择单归属还是多归属，是一个可以精确刻画的决策：\mn{辨析：单归属 vs 多归属刻画的是用户的平台选择行为，与"排他性协议"不同：后者是平台以契约强制单归属（"二选一"），前者是用户的自发选择。}多归属的收益是第二个平台带来的增量网络价值，代价是多付一份接入费。

\begin{derivation}{多归属决策}{}
\textbf{多归属决策的临界条件。}设用户加入平台 $i$ 的效用为 $u_i = v + \theta n_i - p_i$（$v$ 为平台基本价值，$\theta n_i$ 为网络价值），两个平台分别提供 $n_1$、$n_2$ 的网络规模。同时加入两平台的效用为
\[
u_{12} = v + \theta(n_1+n_2) - p_1 - p_2
\]
（多归属可享受两个网络的完整价值，这要求两平台的网络互不重叠；若两平台用户高度重合，增量网络价值小于 $\theta n_2$，多归属的吸引力相应减弱）。只加入平台 $1$ 的效用为 $u_1 = v + \theta n_1 - p_1$。多归属优于单归属的条件为 $u_{12} \ge u_1$，即
\[
\theta n_2 \ge p_2,
\]
对称地，多归属优于只加入平台 $2$ 的条件为 $\theta n_1 \ge p_1$。含义直观：只有当第二个平台的\important{增量网络价值}超过其价格时，用户才愿意多归属。推论：当网络外部性很强（$\theta$ 大）而平台差异化弱（两平台的用户群体高度重合）时，多归属的增量价值小，用户趋于单归属，平台间的竞争焦点转向"争夺独家用户"；当平台差异化强时，多归属普遍，平台竞争趋于缓和（各平台在自己的差异化空间内垄断定价）。这一结论解释了为什么平台竞争最激烈的领域（外卖、网约车）恰恰是差异化最弱的领域：低差异化下多归属的增量网络价值有限，用户对平台没有忠诚，平台只能以补贴争夺用户的每一笔交易。
\end{derivation}

多归属与单归属的区分直接呼应第 3 章赢家通吃的第四个条件（多归属成本高）。现实中两侧的归属模式常常不对称：商家侧普遍多归属（同时入驻多个平台的成本低、收益大），消费者侧则视转换成本而定。这一观察不只是描述性的，它正是下一小节竞争瓶颈模型的现实对应：一侧单归属、一侧多归属的市场结构设定了竞争舞台，而\important{哪一侧更离不开对方，谁就承受高价}，两者共同决定谁承受高价、谁获得补贴。

\subsection{Armstrong 竞争瓶颈模型}

上一小节指出两侧的归属模式常常不对称，这种不对称本身就构成一种市场结构。Armstrong 的竞争瓶颈模型刻画的正是其中最具策略性的一种：一侧用户单归属，另一侧多归属。单归属的一侧成为平台之间争夺的对象，而平台一旦争得这些用户，就对通往他们的接入通道形成独占，这条通道即所谓"瓶颈"。

\begin{definition}[竞争瓶颈（competitive bottleneck）]
竞争瓶颈是指双边市场中一侧用户单归属、另一侧用户多归属时形成的市场结构：平台为争夺单归属侧展开激烈竞争，同时对该侧用户的接入通道形成垄断，多归属侧若要触达全部用户，只能逐个平台分别接入。
\end{definition}

竞争瓶颈下的价格结构无法从单侧的弹性直接读出，两侧的一阶条件互为条件，必须联立求解，下面给出完整推导。

\begin{derivation}{竞争瓶颈模型}{}
设两个平台对称竞争。\textbf{卖方单归属}：卖家只能入驻一个平台，加入平台 $1$ 的效用为
\[
U_S^1 = \alpha_S n_B^1 - p_S^1,
\]
其中 $n_B^i$ 为平台 $i$ 的买方规模，$\alpha_S$ 为每个买方给卖方带来的效用增量（交叉网络外部性），$p_S^i$ 为平台 $i$ 对卖方的收费。由 Hotelling 差异化结构（4.3.3 节详述）导出的单归属卖方需求为
\[
n_S^1 = \frac12 + \frac{\alpha_S(n_B^1-n_B^2) - (p_S^1-p_S^2)}{2t_S},
\]
其中 $t_S$ 为卖方侧差异化参数：卖方对平台的偏好差异越大（$t_S$ 大），单一平台对卖方群的锁定越强；两侧单归属，$n_S^1+n_S^2=1$。\textbf{买方多归属}：买方同时接入两个平台，平台 $i$ 的买方规模为
\[
n_B^i = \phi(\alpha_B n_S^i - p_B^i),
\]
其中 $\phi(\cdot)$ 为递增函数（$\phi' > 0$），$\alpha_B n_S^i$ 为平台 $i$ 卖方规模给买方带来的价值，$p_B^i$ 为买方侧价格；同一买方可同时计入两个平台的规模。平台 $1$ 的利润为
\[
\pi_1 = (p_S^1-f_S)n_S^1 + (p_B^1-f_B)n_B^1
\]
（$f_S, f_B$ 为两侧边际成本）。对 $p_S^1$ 求一阶条件（链式法则，$\frac{\partial n_B^1}{\partial p_S^1} = \alpha_B\phi'\frac{\partial n_S^1}{\partial p_S^1}$）：
\[
n_S^1 - \frac{p_S^1-f_S}{2t_S} - \frac{\alpha_B\phi'(p_B^1-f_B)}{2t_S} = 0,
\]
即
\[
p_S^1-f_S = 2t_S n_S^1 - \alpha_B\phi'(p_B^1-f_B).
\]
对 $p_B^1$ 求一阶条件（对称点处对手平台的买方规模不随本平台价格变动，$\frac{\partial n_S^1}{\partial p_B^1} = -\frac{\alpha_S\phi'}{2t_S}$）：
\[
n_B^1 - \phi'(p_B^1-f_B) - \frac{\alpha_S\phi'}{2t_S}(p_S^1-f_S) = 0,
\]
即
\[
p_B^1-f_B = \frac{n_B^1}{\phi'} - \frac{\alpha_S}{2t_S}(p_S^1-f_S).
\]
对称均衡中 $n_S^1 = 1/2$、$n_B = \phi(\alpha_B/2 - p_B)$，两个一阶条件给出均衡价格结构：
\[
p_S - f_S = t_S - \alpha_B\phi'(p_B-f_B), \qquad
p_B - f_B = \frac{n_B}{\phi'} - \frac{\alpha_S}{2t_S}(p_S-f_S).
\]
这两个一阶条件互为反馈，必须联立求解才能确定两侧的绝对价格水平。记单归属侧加成为 $x = p_S^1 - f_S$、多归属侧加成为 $y = p_B^1 - f_B$。将第一式代入第二式消去 $x$，整理得
\[
\left(1 - \frac{\alpha_S\alpha_B\phi'}{2t_S}\right) y = \frac{n_B}{\phi'} - \frac{\alpha_S}{2},
\]
其中括号项刻画两侧交叉外部性的自我强化：买方增加抬高卖方对平台的估值，卖方增加又抬高买方的估值，反馈沿两侧交替放大。为保证价格反馈不发散，必须有 $1 - \alpha_S\alpha_B\phi'/(2t_S) > 0$，即稳定性条件
\[
\alpha_S\alpha_B\phi' < 2t_S.
\]
在此条件下解得两侧均衡加成的显式解：
\[
y = \frac{n_B/\phi' - \alpha_S/2}{1 - \alpha_S\alpha_B\phi'/(2t_S)}, \qquad
x = t_S - \alpha_B\phi' y.
\]
分母恒正，$y$ 的符号完全由分子决定，均衡价格结构按 $\alpha_S\phi'$ 与 $2n_B$ 的大小关系分为两种情形。情形一：$\alpha_S\phi' < 2n_B$，此时 $y > 0$，代入得 $x < t_S$，即多归属侧支付正的弹性加成，单归属侧的价格反而低于 Hotelling 垄断加成，平台把多归属侧的利润让渡给了单归属侧。情形二：$\alpha_S\phi' > 2n_B$，此时 $y < 0$，代入得 $x > t_S$，即多归属侧获得净补贴，补贴成本由单归属侧高于垄断加成的超额收费承担。分界线 $\alpha_S\phi' = 2n_B$ 处多归属侧恰好按边际成本定价。竞争瓶颈的结论因此取决于两侧参数：\important{当单归属侧对多归属侧的依赖较弱时，多归属侧付垄断加成、单归属侧获让利；当单归属侧极度依赖多归属侧导流时，价格结构翻转，多归属侧获补贴、单归属侧付高价}。
\end{derivation}

情形一正是 Armstrong 原模型的设定及其经典结论：报纸的读者单归属，但读者对单一报纸的依赖有限，报社以低价甚至免费争夺读者，再向必须逐家投放才能覆盖全部读者的广告主收取高价；游戏主机同理，主机降价争夺玩家，再向多归属的游戏开发者抽成。此时平台对多归属侧握有接入垄断权，多归属侧成为被榨取的一方，这一方向是教材与考试的标准答案。\mn{竞争瓶颈标准结论（Armstrong 方向）：多归属侧付垄断加成、单归属侧获让利，如报纸补贴读者、向广告主收费。判断线 $\alpha_S\phi' = 2n_B$：单归属侧对多归属侧的依赖一旦超过多归属侧的自身弹性项，结构翻转，多归属侧反获补贴。}情形二则对应中国平台经济的现实：电商消费者普遍多归属、随时比价，而商家被"二选一"协议锁在单一平台上，销量高度依赖平台导流，$\alpha_S$ 极大，价格结构于是翻转，消费者拿补贴、商家承担高额佣金。两种情形背后是同一规律：\important{哪一侧更离不开对方，哪一侧付高价}。

这一规律给出了"二选一"禁令的经济学机理：禁止"二选一"恢复了商家的多归属自由，竞争瓶颈结构随之被拆解，商家侧高于垄断加成的超额收费也就失去了支撑。中国的监管实践印证了这一机理：《关于平台经济领域的反垄断指南》（国务院反垄断委员会，2021）明确将"二选一"列为可能构成滥用市场支配地位的限定交易行为；同年，市场监管总局对阿里巴巴集团的"二选一"行为作出行政处罚，罚款 182.28 亿元。第 6 章将展开该案的认定框架与救济措施。

\subsection{平台竞争的 Hotelling 扩展}

多归属与竞争瓶颈分别刻画了平台竞争的两种典型结构（全部单归属与一侧单归属），本节考察对称的基准情形：\textbf{两侧用户均单归属}。此时平台竞争由 Hotelling 差异化模型刻画：消费者沿长度为 $1$ 的线段均匀分布，两个平台位于两端，用户到平台的"距离成本"代表对产品的偏好差异，距离成本参数 $t$ 越大差异化越强、竞争越缓和（4.4 节补贴战模型将再次调用这一结构）。

\begin{derivation}{Hotelling 模型的双边扩展}{}
设两侧用户均单归属、各自在线段上均匀分布。侧 $k$（$k \in \{B,S\}$）用户在位置 $x$ 处加入平台 $1$ 的效用为
\[
u_k^1 = \alpha_k n_j^1 - p_k^1 - t_k x, \qquad j \ne k,
\]
其中 $\alpha_k$ 为另一侧每个用户给侧 $k$ 用户带来的效用增量（交叉网络外部性），$t_k$ 为侧 $k$ 的差异化参数，$n_j^i$ 为平台 $i$ 另一侧的规模。边际用户满足 $u_k^1 = u_k^2$，结合理性预期 $n_k^1 = x$、$n_k^2 = 1 - n_k^1$，得两侧需求
\[
n_k^1 = \frac12 + \frac{\alpha_k(n_j^1-n_j^2) + (p_k^2-p_k^1)}{2t_k}.
\]
平台 $1$ 的利润 $\pi^1 = (p_B^1-c_B)n_B^1 + (p_S^1-c_S)n_S^1$（$c_k$ 为侧 $k$ 的边际成本）。设 $t_Bt_S > \alpha_B\alpha_S$（差异化足够强、交叉外部性不过大，保证需求良定义）。在对称点对 $p_B^1$ 求一阶条件，先由两侧需求方程对 $p_B^1$ 全微分：
\[
\frac{\partial n_B^1}{\partial p_B^1} = -\frac{1}{2t_B} + \frac{\alpha_B}{t_B}\frac{\partial n_S^1}{\partial p_B^1}, \qquad
\frac{\partial n_S^1}{\partial p_B^1} = \frac{\alpha_S}{t_S}\frac{\partial n_B^1}{\partial p_B^1},
\]
消元得
\[
\frac{\partial n_B^1}{\partial p_B^1} = -\frac{t_S}{2(t_Bt_S-\alpha_B\alpha_S)}, \qquad
\frac{\partial n_S^1}{\partial p_B^1} = -\frac{\alpha_S}{2(t_Bt_S-\alpha_B\alpha_S)}.
\]
代入 $\frac{\partial\pi^1}{\partial p_B^1} = 0$（对称点处 $n_B^1 = 1/2$）：
\[
\frac12 - \frac{(p_B-c_B)t_S}{2(t_Bt_S-\alpha_B\alpha_S)} - \frac{(p_S-c_S)\alpha_S}{2(t_Bt_S-\alpha_B\alpha_S)} = 0,
\]
乘 $2(t_Bt_S-\alpha_B\alpha_S)$ 得
\[
(p_B-c_B)t_S + (p_S-c_S)\alpha_S = t_Bt_S-\alpha_B\alpha_S.
\]
对 $p_S^1$ 求导对称地给出 $(p_S-c_S)t_B + (p_B-c_B)\alpha_B = t_Bt_S-\alpha_B\alpha_S$。两式联立解得对称均衡价格（代入验证：$(t_B-\alpha_S)t_S + (t_S-\alpha_B)\alpha_S = t_Bt_S-\alpha_B\alpha_S$，第二式同理）：
\[
p_B = c_B + t_B - \alpha_S, \qquad p_S = c_S + t_S - \alpha_B.
\]
\end{derivation}

均衡价格的三个组成部分各有含义：边际成本 $c_k$、差异化加成 $t_k$、交叉外部性折扣 $\alpha_j$（$j$ 为另一侧）。与单侧 Hotelling 均衡 $p = c+t$ 相比，双边化使每侧价格下降 $\alpha_j$：该侧用户为另一侧创造的外部性价值，平台以降价的形式让渡给该侧。\mn{Hotelling 双边扩展的价格公式：$p_k = c_k + t_k - \alpha_j$（$j$ 为另一侧）。每侧价格在标准 Hotelling 水平上扣减该侧用户对另一侧的外部性价值。}由此，差异化弱的一侧（$t_k$ 小）竞争激烈、价格趋近边际成本；当 $\alpha_j > t_k$ 时该侧价格低于边际成本，$\alpha_j > t_k + c_k$ 时为负（补贴）。综合本节三个模型，双边市场竞争分析的统一逻辑是：\important{两侧竞争结构对称时，竞争约束强的一侧（低差异化、低转换成本）价格低}；两侧结构不对称时（竞争瓶颈），价格方向还要看哪一侧更依赖对方，多归属侧可能反而付垄断加成。复习时以"对称看竞争、不对称看依赖"八个字串起三个模型。

并非所有双边市场都走向一家独大，第 3 章赢家通吃的条件在这里有一个重要补充：\textbf{负向网络外部性}（negative network externality）对平台规模的制约。交叉外部性未必恒为正。其一，\textbf{同边负外部性}：同侧用户之间的竞争损害（外卖平台上同商圈商家的价格竞争、招聘平台上同岗位求职者的相互竞争），用户规模越大，同侧参与者的人均匹配机会反而下降；其二，\textbf{交叉负外部性}：一侧规模过度扩张降低另一侧的匹配质量，商家过多使消费者的搜寻成本上升，海量低质量供给淹没高质量供给（逆向选择），求职者过多使雇主筛选成本上升。\mn{负向网络外部性的两个来源：同边竞争（求职者互相竞争）与交叉拥堵（商家过多、搜寻成本上升）。典型平台：招聘与婚恋平台因匹配质量随规模恶化而不呈现赢家通吃，多平台长期共存。}负向外部性的存在解释了招聘平台（前程无忧、BOSS 直聘、猎聘）与婚恋平台长期多强并存的现实：这类平台的匹配质量依赖筛选而非规模，规模扩大带来的噪音与拥堵抵消了网络价值，赢家通吃的力量被负向外部性系统性削弱。平台的竞争战略因此不仅是"做大"，还包括"过滤"：审核机制、信用评分与算法排序的功能，正是把负向外部性转化为平台的组织能力，第 12 章的平台责任制度讨论即由此展开。

\section{平台策略行为}

\subsection{排他性协议}

平台竞争中最常见的契约策略是排他性协议：平台要求一侧用户只在本平台交易，俗称"二选一"。其竞争效应的分析框架是 Armstrong 竞争瓶颈的延伸：排他性协议以契约把商家从多归属强制变为单归属，同时切断商家在对手平台获取流量的通道，使其对缔约平台的依赖加深，价格结构翻转为商家侧承受超额收费（4.3.2 节的情形二）。

\begin{definition}[排他性协议（exclusive dealing）]
排他性协议是平台与一侧用户签订的契约安排，要求该用户只在本平台交易、不得加入竞争平台，俗称"二选一"。
\end{definition}

排他性协议之所以值得单独分析，在于它能够阻止竞争者进入。下面用一个市场圈定的简化模型给出进入阻止条件，并把双边市场与单边市场的圈定效果作一对照。

\begin{derivation}{市场圈定与进入阻止条件}{}
设在位平台 $I$ 已在市场中运营，潜在进入者 $E$ 考虑是否进入。卖方总量标准化为 $1$，在位平台与其中比例 $\lambda \in [0,1]$ 的卖方签订排他性协议，这部分卖方不得加入进入者的平台。进入需支付固定成本 $F$（技术开发、市场培育、初期补贴），记进入者在不受圈定（$\lambda = 0$）时可获得的毛利润为 $\pi_0 > 0$。

第一步，确定圈定对进入者卖方规模的影响。签约卖方被锁定后，进入者只能接触剩余卖方，其卖方侧规模为
\[
n_S^E = 1 - \lambda.
\]

第二步，把交叉网络外部性引入买方侧。买方接入进入者平台的价值随该平台的卖方规模递增；在买方价值与卖方规模成比例的设定下，进入者的买方侧规模同比例缩减：
\[
n_B^E = (1-\lambda)\,\bar n_B,
\]
其中 $\bar n_B$ 为不受圈定时进入者可获得的买方规模。

第三步，写出进入者的毛利润。沿用 4.2 节 $V = D_BD_S$ 的交易量设定，交易量正比于两侧规模之积，故
\[
\pi_E(\lambda) = (1-\lambda)^2\,\pi_0 .
\]
\important{排他协议对进入者利润的削减是平方的}：卖方侧被圈定一次，买方侧再通过交叉外部性被削减一次。

第四步，求进入阻止条件。进入者只在毛利润足以覆盖固定成本时进入，故进入被阻止当且仅当 $\pi_E(\lambda) < F$，即
\[
(1-\lambda)^2\pi_0 < F
\;\Longrightarrow\;
1-\lambda < \sqrt{F/\pi_0}
\;\Longrightarrow\;
\lambda > 1 - \sqrt{F/\pi_0}.
\]
记临界锁定比例
\[
\lambda^* = 1 - \sqrt{F/\pi_0},
\]
在位平台的签约覆盖率只要超过 $\lambda^*$，进入即被阻止。

第五步，与单边市场对照。若不存在交叉网络外部性，进入者的利润仅随卖方规模线性缩减，$\pi_E = (1-\lambda)\pi_0$，进入阻止条件为 $\lambda > 1 - F/\pi_0$，临界值 $\lambda^*_{\text{单边}} = 1 - F/\pi_0$。记 $r = F/\pi_0 \in (0,1)$，由 $\sqrt r > r$ 得
\[
\lambda^* = 1-\sqrt r \;<\; 1-r = \lambda^*_{\text{单边}} .
\]
\important{双边市场阻止进入所需的排他覆盖比例低于单边市场}。以 $r = 0.25$ 为例，双边市场锁定一半卖方（$\lambda^* = 0.5$）即可阻止进入，单边市场则需锁定四分之三（$\lambda^*_{\text{单边}} = 0.75$）。
\end{derivation}

这一条件揭示了排他性协议的竞争损害机制：协议压缩进入者可触达的两侧规模，使潜在竞争者的利润无法覆盖进入成本，进入因此被阻止，在位平台得以维持 4.3.2 节情形二的价格结构，继续向被锁定的商家侧收取超额费用。\mn{圈定的平方效应速记：卖方被锁一次、买方随之少一次，$(1-\lambda)^2$。临界锁定率 $\lambda^*=1-\sqrt{F/\pi_0}$，比单边市场的 $1-F/\pi_0$ 更低，故"二选一"在平台经济中危害更大。}平方项是双边市场特有的放大器，它意味着排他协议只需覆盖相对小的一部分卖方就足以关上进入的大门，这正是"二选一"在平台经济中比在传统零售中危害更大的根本原因。

排他性协议同时具有效率促进的一面，福利评价需要权衡两种效应。\textbf{竞争损害效应}如上所述，包括进入被阻止、被锁定商家承受超额收费，以及对手平台因失去商家资源而在买方侧竞争减弱。\textbf{效率促进效应}来自专属投入的激励：平台若为商家提供培训、流量扶持或定制化系统，而商家转身把这些投入带往对手平台使用，平台的投资激励就会被削弱，排他性协议通过防止这种"搭便车"维持了投资激励，平台间的竞争也在协议期内从价格竞争转向服务竞争。净效应取决于两者的相对大小：当平台市场势力强、协议覆盖比例超过 $\lambda^*$、又缺乏可验证的专属投入时，损害效应占优，排他性协议构成垄断行为。

现实中排他性协议以多种面貌出现：独家经营奖励（平台向承诺独家入驻的商家提供流量与费率优惠）、独家内容版权，以及强制性的"二选一"。前两类附有对价、伴随平台专属投入，属于普遍存在的商业实践；强制"二选一"缺乏对价与效率理由，成为执法重点（第 6 章阿里巴巴案）。定性判断的关键不在协议形式，而在 Armstrong 框架给出的两个变量：\important{平台市场势力的大小，以及排他是否实质剥夺了用户的多归属选择}。

排他性协议之外，平台对商家施加的契约约束还有两类常见形式，三者同属\textbf{纵向约束}（vertical restraints）的范畴。其一是\textbf{价格承诺}，最典型的是\textbf{最惠国待遇条款}（most-favored-nation clause，MFN），平台要求商家在本平台的报价不高于其在其他渠道的报价。这类条款表面上保护消费者，实际效果却是消解商家以低价吸引流量的动机：商家无法在新平台上给出更低价格，进入者也就失去了最有效的获客手段，圈定效果与排他性协议相似，手段则从"禁止加入"变成"禁止低价"。酒店预订平台与电商平台的价格平价条款已在欧洲多国受到反垄断审查。其二是地域限制、客户限制与最低转售价格维持等传统纵向约束，平台以交易规则的形式限定商家的经营范围与定价自由。评价这些约束的标准与排他性协议一致：看它们是否实质压缩了竞争平台可触达的用户规模，以及是否伴随可验证的专属投入。中国《电子商务法》（2019 年 1 月 1 日施行）第三十五条对平台利用服务协议、交易规则及技术手段对商家施加不合理限制的行为作出禁止性规定，为这类约束提供了平台经济领域的专门规范。

\subsection{平台包抄}

排他性协议争夺的是同一边的用户，平台竞争的另一种超常规形式则直接跨越市场边界，这就是\textbf{平台包抄}。

\begin{definition}[平台包抄（platform envelopment）]
平台包抄是指一个平台利用其在某一市场的用户基础，通过捆绑相邻功能进入另一平台市场，吸收乃至取代该市场在位平台的竞争策略。
\end{definition}

典型案例：谷歌利用其庞大的搜索用户基础推出谷歌地图，随后借助安卓系统的用户基础推出谷歌钱包，先后包抄独立的地图与支付平台；微信利用社交用户基础推出支付功能，包抄支付宝的支付市场。\mn{案例：平台包抄成立的条件：相邻市场与主市场共享用户基础 + 交叉网络外部性可被复用 + 被包抄方缺乏独特资产。}包抄策略的经济学本质是用户基础与交叉网络外部性的跨市场复用，其不对称优势可以形式化。

\begin{derivation}{包抄的成本与规模优势}{}
设主市场（市场 A）的平台拥有用户基础 $N_A$，相邻市场（市场 B）的在位者拥有用户基础 $N_B$。记\textbf{用户重叠度} $\rho \in [0,1]$ 为主市场用户中对市场 B 的服务同样存在需求的比例，则包抄者可直接转化的用户规模为 $\rho N_A$。在市场 B 中每获取一个用户的成本记为 $c_B$（广告费、补贴等）。

第一步，比较两者的获客成本。在位者获取 $n$ 个新用户需支出 $c_Bn$；包抄者通过主产品的功能捆绑（预装、默认入口、消息推送）把市场 B 的服务分发给既有用户，边际分发成本为 $\epsilon \approx 0$，获取同样 $n$ 个用户仅需支出 $\epsilon n$。包抄者的成本优势为
\[
\Delta = (c_B-\epsilon)\,n \approx c_B\,n,
\]
即在位者的全部获客支出构成包抄者的净优势。

第二步，比较进入时的初始规模。包抄者一进入市场 B 即拥有规模 $\rho N_A$，它超过在位者当且仅当
\[
\rho N_A > N_B
\;\Longleftrightarrow\;
\rho > \frac{N_B}{N_A}.
\]
\important{包抄成立的规模条件是用户重叠度超过两个市场的规模比}。$N_A \gg N_B$ 时该条件极易满足：主市场的用户基础越大，包抄所需的重叠度越低。

第三步，引入市场 B 的网络外部性。设市场 B 的临界规模为 $n^c$（第 3 章）。包抄者若满足
\[
\rho N_A > n^c,
\]
则进入的瞬间即越过临界规模，无须经历 4.1.3 节所述的启动阶段，鸡生蛋问题对包抄者不构成约束。合并两个条件，包抄成立要求
\[
\rho N_A > \max\{N_B,\; n^c\}.
\]
该式表明包抄的可行性由\important{用户重叠度与主市场规模的乘积}决定：$\rho \to 0$ 时，无论 $N_A$ 多大，包抄都不成立。
\end{derivation}

上述条件点明了包抄的经济学本质：包抄依靠用户基础与交叉网络外部性的跨市场复用，而复用的前提是两个市场的用户群重合。若两市场用户互不相干，例如谷歌搜索用户与专业绘图软件用户，$\rho$ 接近零，包抄优势不复存在。

对反垄断而言，平台包抄使"市场进入"的常规含义失效。传统的进入壁垒分析考察同一市场内新企业的进入难度，而包抄者来自相邻市场，既不需要从零积累用户，也不需要跨越临界规模，常规意义上的进入壁垒对它形同虚设。进入壁垒的分析因此必须扩展到相邻市场，考察哪些巨头的用户群覆盖了本市场，第 6 章的相关市场界定与市场势力测度正是沿这一方向对传统框架的修正。

包抄对在位者的含义是：\important{市场的边界由用户基础定义，而非由产品定义}，只要用户群重合，巨头就随时可能跨界进入。微信支付之于支付宝即是明证：独立平台即使在本市场建立了很深的壁垒，只要相邻巨头的用户群覆盖自己的用户群、且交叉外部性可以迁移，其市场地位就始终面临被包抄的可能。用户群不重合的专业市场（如工业软件）则天然免疫于消费平台巨头的包抄。

\subsection{案例：外卖平台补贴战的经济学}

外卖平台的竞争史是本章理论的综合演练，倾斜式定价、竞争瓶颈与临界规模三条线索在这里同时起作用。补贴战的动力学可以用一个简化的双平台模型完整刻画，模型把霍特林差异化结构与交叉网络外部性结合起来，回答三个问题：补贴为什么有效、补贴为什么会升级、补贴什么时候停止。

\begin{derivation}{补贴战的动态学}{}
考虑两个对称平台 $A$、$B$ 竞争同一市场。消费者沿长度为 $1$ 的霍特林线段均匀分布，平台位于两端；消费者加入平台 $i$ 的效用为
\[
u_i = v + \beta n_i - p_i - \tau x,
\]
其中 $\beta n_i$ 为交叉网络外部性价值（$n_i$ 为平台 $i$ 的用户份额），$p_i$ 为消费者侧价格（补贴时 $p_i < 0$），$\tau$ 为单位距离成本（差异化程度）。

第一步，求用户份额。边际消费者满足 $u_A = u_B$；该消费者左侧的区间长度即平台 $A$ 的用户份额，故其位置 $x = n_A$，右侧区间长度 $1-n_A$ 为平台 $B$ 的份额。代入理性预期（实际份额与预期一致），边际条件为
\[
v + \beta n_A - p_A - \tau n_A = v + \beta n_B - p_B - \tau(1 - n_A),
\]
利用 $n_A + n_B = 1$ 解得平台 $A$ 的用户份额：
\[
n_A = \frac12 + \frac{p_B - p_A}{2(\tau - \beta)}.
\]

第二步，求价格优势的放大倍数。对上式关于 $p_A$ 求导，得
\[
\left|\frac{\partial n_A}{\partial p_A}\right| = \frac{1}{2(\tau-\beta)},
\]
即降价一单位使份额扩大 $1/[2(\tau-\beta)]$。$\beta$ 趋近 $\tau$ 时该倍数趋于无穷；$\beta \geq \tau$ 时分母非正，份额函数不再给出内点解，\important{市场直接倾覆}（赢家通吃）。

第三步，求不跟进降价的利润损失。若平台 $B$ 降价 $\Delta$ 而平台 $A$ 不跟进，$A$ 的份额损失为 $\delta n = \Delta/[2(\tau-\beta)]$。平台 $A$ 的利润同时来自两侧：$\pi_A = p_A n_A + \Pi^S(n_A)$，其中 $\Pi^S(n_A)$ 为商家侧佣金利润，随消费者规模递增（$\Pi^S{}' > 0$：买家越多，商家越愿意入驻付费）。份额下降 $\delta n$ 带来的利润损失为
\[
\bigl(\Pi^S{}'(n_A) - |p_A|\bigr)\cdot\frac{\Delta}{2(\tau-\beta)},
\]
即佣金损失 $\Pi^S{}'\delta n$ 减去节省下来的补贴支出 $|p_A|\delta n$。当佣金敏感度超过补贴强度时该损失为正，且随 $\beta$ 趋近 $\tau$ 被无限放大。

第四步，求停战条件。平台 $i$ 参与补贴战的理性条件是预期收益的净现值非负：
\[
PV_i = -S_i + \sum_{t=T}^{\infty}\delta^t\,\pi^M_i = -S_i + \frac{\delta^T}{1-\delta}\,\pi^M_i \geq 0,
\]
其中 $S_i$ 为从现在到对手退出的预期补贴投入现值，$T$ 为对手退出时刻，$\pi^M_i$ 为赢家通吃后的垄断佣金利润，$\delta \in (0,1)$ 为贴现因子（反映融资成本与耐心），末式由等比级数求和化简得到。该式给出\important{补贴战强度的三个决定变量}：赢家通吃的速度（$\beta$ 与 $\tau$ 的差距决定 $T$）、垄断利润的规模（$\pi^M$）、以及耐心程度（$\delta$，取决于资本市场的融资条件）。当 $PV_i < 0$ 对双方同时成立时，默契停战成为均衡，市场停留于双寡头。
\end{derivation}

四个步骤的结果合起来解释了补贴战的行为逻辑。放大倍数说明补贴为什么有效：网络外部性越强，同样幅度的降价买到的份额越多，平台因此愿意为看似微小的价格优势投入巨额资金。利润损失公式说明补贴为什么会升级：不跟进的一方承受的佣金损失随外部性强度放大，跟进降价成为占优选择，双方交替降价，补贴竞赛由此形成。停战条件则说明补贴为什么会停止：当双方都判断击垮对手所需的投入超过垄断阶段的预期回报时，继续投入不再理性。2014 年爆发的滴滴快的补贴战于 2015 年 2 月以合并告终，正是因为双方在势均力敌的投入强度下都无法形成单方胜利的预期，补贴规模逼近了停战条件的边界。

这一模型同时解释了补贴战的现实细节。美团与饿了么的长期对峙中，补贴呈现"脉冲式"特征，节假日与新市场开拓时骤然升级，随后退潮。脉冲对应于双方对临界规模的反复试探：每次补贴升级都试图把对手压制到其临界规模之下，使其网络收缩消亡，或者把己方网络推过临界规模（第 3 章）；试探未能达到目的时，退守止损是理性选择。外卖平台同时连接消费者（价格敏感、多归属、交叉外部性贡献大）与商家（单归属倾向、依赖平台流量），按 Rochet–Tirole 定价规则向消费者侧倾斜定价、向商家侧收取佣金，又按竞争瓶颈逻辑使单归属的商家侧承受高抽佣。补贴的资金来源正是商家佣金，\important{补贴战的实质是以商家侧的垄断租金资助消费者侧的竞争投入}。平台补贴因此是网络外部性市场中的理性投资，其规模与期限由赢家通吃的预期收益贴现决定。

\section{平台生态系统}

前四节把平台当作连接两侧用户的组织来分析，定价、竞争与策略行为都在"两侧"这一框架内展开。现实中的头部平台早已越出这一框架：阿里巴巴同时经营电商、支付、物流与云计算，腾讯的社交、游戏、支付与内容互相导流。这些平台连接的用户群体远不止两类，各业务之间的外部性也超出了简单的双向关系。本节考察平台形态的这一演进，以及演进所带来的治理问题与模式选择。

\subsection{平台形态的三阶段演进}

平台并非一开始就具备今天的复杂形态。从连接两类用户的双边平台，到连接多类用户的多边平台，再到以平台为核心的生态圈，平台形态经历了三个阶段的演进。三个阶段的划分依据与转变动力是名词解释与简答题的常考内容，需要能够准确复述。

\begin{definition}[多边平台（multi-sided platform）]
多边平台是连接三类及以上相互依赖的用户群体的平台组织：各类用户群体之间存在多向的交叉网络外部性，平台需要同时管理多组价格结构。
\end{definition}

\begin{definition}[平台生态圈（platform ecosystem）]
平台生态圈是以核心平台为中心、由各类互补性参与者共同构成的动态创新联合体：参与者围绕平台提供的基础设施与规则展开互补性创新，平台则通过对接口、数据与规则的控制，实现交叉网络外部性的最大限度内部化。
\end{definition}

三个阶段的\textbf{划分依据}有三个维度。其一，\important{网络外部性的内部化空间}：双边平台只能内部化两类用户之间的一组交叉外部性，多边平台可以内部化多组，生态圈则借助互补品之间的相互增强，把外部性的内部化推到最大限度。其二，\important{业务多元化程度}：双边平台聚焦单一业务场景，多边平台在相关业务上横向扩展，生态圈覆盖跨行业的多元业务组合。其三，\important{赋能能力}：双边平台提供撮合，多边平台提供撮合与基础服务，生态圈则向参与者输出技术能力、数据能力与信用能力，参与者的经营活动在相当程度上依托平台的赋能。\mn{三形态速记：双边（两类用户、一组外部性）→ 多边（多类用户、多组外部性）→ 生态圈（互补者联合体、外部性内部化最大化）。划分三维：内部化空间、多元化程度、赋能能力；转变两动力：协同性提升、动态竞争压力。}三个维度沿同一方向递进，构成形态演进的判别标准。

演进的\textbf{转变动力}有两个。其一，\important{业务协同性的提升}：随着业务范围扩大，各业务在用户、数据与基础设施上的共享程度提高，协同收益超过管理成本，平台因而有动力把边界继续外推。其二，\important{技术高速迭代下的动态竞争压力}：数字技术的迭代速度使单一业务的领先地位难以长期维持，4.4.2 节的包抄分析已经说明，只要用户群重合，相邻市场的巨头随时可能跨界进入。构建生态圈的防御含义正在于此：业务组合越丰富、参与者与平台的绑定越深，被单点包抄的风险越低。

生态圈对参与者形成的锁定与单一平台的锁定有性质上的差别，可称为\important{多维交叉锁定}：用户在电商、支付、物流与内容多个维度同时与平台建立关系，任何一个维度的转换成本单独看都不算高，但多个维度的转换必须同步完成，合并起来的转换成本远超各维度之和。第 3 章讨论的锁定效应在生态圈中因此被显著强化，这也是生态型平台的市场地位格外稳固的原因。

阿里巴巴的演进提供了三阶段的完整样本。1999 年成立的阿里巴巴以 B2B 信息撮合起步，2003 年推出的淘宝是典型的双边平台，连接买家与卖家。2004 年支付宝独立运营，2009 年阿里云成立，2013 年菜鸟网络组建，平台连接的群体扩展到支付机构、开发者、物流服务商与金融机构，形态转入多边平台。此后，商家在阿里体系内可以同时获得交易撮合、支付结算、物流履约、云计算资源与信贷融资，各项服务相互调用、数据相互沉淀，生态圈形态由此成型。腾讯以社交为核心，通过微信开放平台把公众号、小程序与支付能力提供给外部开发者，走的是同一条路径。

\subsection{平台治理}

生态圈的参与者众多、利益诉求各异，平台需要一套规则体系来协调，这套规则体系即平台治理。平台治理既决定生态圈能否持续扩张，也决定平台与参与者之间的利益分配。它是近年出现频率上升较快的考点，三项内容需要能够分条作答。

\begin{definition}[平台治理（platform governance）]
平台治理是平台对生态系统中的参与者、交易与互动所施加的规则安排：平台通过准入规则、交易规则、争议解决机制与技术接口设计，规定谁可以参与、以何种方式参与、利益如何分配。
\end{definition}

Parker 与 Van Alstyne 把平台治理的内容概括为三个方面。其一，\important{界定平台的边界与开放性程度}：平台需要决定哪些功能由自己提供、哪些交给第三方，以及向第三方开放到什么程度。开放度过低会抑制互补者的创新积极性，生态圈难以扩张；开放度过高则平台难以保证服务质量，核心业务也容易被互补者反向侵蚀。iOS 与 Android 在应用生态上的不同选择正是开放度决策的两个方向。\mn{平台治理三内容速记："定边界、调关系、平利益"，即界定平台边界与开放性程度、调和内外部关系、平衡各类参与者利益（Parker–Van Alstyne 框架）。答题时每条各展开一句即可成段。}其二，\important{调和内外部关系}：平台的自营业务与第三方参与者常常在同一市场竞争，平台既是规则制定者又是市场参与者，这一双重身份要求它在流量分配、搜索排序与费率设定上对内外部主体保持一致，否则会引发自我优待的争议（第 6 章详述）。其三，\important{平衡各类参与者的利益}：生态圈中的买方、卖方、开发者与服务商诉求各异，平台需要在费率、规则与资源分配上取得平衡，任何一方的利益被过度挤压，都会削弱生态圈的长期价值。

平台治理与 4.3.3 节讨论的负向网络外部性直接相关。规模扩大带来的拥堵与匹配质量下降，本质上要依靠治理手段来抑制：审核机制过滤低质量供给，信用评分区分参与者质量，算法排序决定注意力的分配。治理能力因而构成平台的核心竞争力之一，它把负向外部性转化为可管理的组织问题。中国的监管实践也在向这一方向施加压力：《电子商务法》对平台的资质审核、信息公示与安全保障义务作出规定，把部分治理内容上升为法定责任，平台责任制度的完整讨论见第 12 章。

\subsection{平台模式与自营模式}

平台治理的第一项内容是界定边界，其中最基础的一个决策是：平台自己下场经营，还是把经营交给第三方？这一决策区分了平台模式与自营模式，也是理解亚马逊、京东等企业混合形态的关键。两种模式的差别可以沿两个维度刻画。

第一个维度是\important{开放性}。平台模式向第三方开放交易机会，商品的种类与供给规模由市场决定，平台收入来自佣金与增值服务；自营模式由企业自行采购与销售，商品结构由企业决定，收入来自进销差价。开放性的经济含义落在供给侧的规模扩张能力上：平台模式的商品丰富度随入驻商家数量增长，自营模式则受限于企业自身的采购与仓储能力。

第二个维度是\important{资产轻重与控制权}，二者互为代价。自营模式需要投入库存、仓储与配送等重资产，承担存货风险与资金占用，换来的是对商品质量、定价与交付体验的完整控制权；平台模式以轻资产运营，规模扩张不受自有资产约束，代价是把商品质量与服务水平的控制权部分让渡给第三方商家，需要以治理规则来弥补。\mn{平台 vs 自营两维速记：开放性（向第三方开放 vs 自行经营）、资产轻重与控制权（轻资产弱控制 vs 重资产强控制）。一句话权衡："开放换规模，自营换品质"。}两个维度合起来给出一条清晰的权衡：\important{开放换规模，自营换品质}。

现实中的头部电商多采用混合模式，在两个维度上取折中。亚马逊以自营起家，随后开放第三方卖家平台，同时把自营积累的仓配能力以物流服务的形式开放给第三方，兼得规模与品质控制。京东的路径类似，早期以自营 3C 品类与自建物流建立品质与时效优势，此后逐步引入第三方商家扩展品类，自营与平台业务并行。淘宝则始终以平台模式为主，商品丰富度居前，品质控制主要依靠评价体系、信用机制与平台规则实现。三者的差异说明模式选择取决于品类特征与竞争定位：标准化程度高、品质风险大、时效敏感的品类适合自营，长尾品类与个性化需求则适合平台模式。这一权衡在第 11 章讨论产业数字化时还会以企业边界问题的形式再次出现。

\begin{chapterbox}
\textbf{本章考点清单}

\begin{enumerate}
\item 双边市场的判定标准：价格结构非中性 \important{【专硕·核心】}
\item 交叉网络外部性与鸡生蛋问题：零均衡的自我实现、临界规模门槛 $\underline{n}_B = v_S^{-1}(p_S)$ 与三种破解策略 \important{【专硕·核心】}
\item Rochet–Tirole 模型：$V = D_BD_S$ 设定、价格结构 $\frac{p_B}{\eta_B} = \frac{p_S}{\eta_S}$ 与总价格水平 $M = \frac{\eta_B+\eta_S}{\eta_B+\eta_S-1}c$ 的推导 \important{【专硕·核心】}
\item 含便利收益的变体：$c$ 处处替换为 $c-b$，$p_B = \frac{\eta_B}{\eta_B-1}(c-b-p_S)$ 与买方侧补贴条件 $p_S > c-b$；银行卡交换费结构 \important{【专硕·核心】}
\item 三种定价目标（利润/交易量/社会最优）的比较：社会最优总价格低于边际成本 \important{【专硕·核心】}
\item 倾斜式定价：补贴哪一侧的规则（弹性 + 外部性贡献）；与掠夺性定价的辨析 \important{【专硕·核心】}
\item 多归属决策条件 $\theta n_2 \geq p_2$ \important{【专硕·核心】}
\item Armstrong 竞争瓶颈：定义、联立一阶条件的求解、$\alpha_S\phi' = 2n_B$ 判断线与两种价格结构（多归属侧付垄断加成；依赖极强时翻转）的推导 \important{【专硕·核心】}
\item Hotelling 双边扩展：对称均衡 $p_k = c_k + t_k - \alpha_j$ 的推导 \important{【专硕·扩展】}
\item 负向网络外部性（同边竞争、交叉拥堵）对平台规模的制约与招聘/婚恋平台的多强并存 \important{【专硕·扩展】}
\item 排他性协议：市场圈定的进入阻止条件 $\lambda^* = 1-\sqrt{F/\pi_0}$、圈定的平方效应，以及竞争损害与效率促进的权衡 \important{【专硕·核心】}
\item 价格承诺（最惠国待遇条款）与纵向约束的圈定效应 \important{【专硕·扩展】}
\item 平台包抄：成立条件 $\rho N_A > \max\{N_B,\,n^c\}$ 与经典案例 \important{【专硕·扩展】}
\item 外卖补贴战案例：倾斜定价、竞争瓶颈与临界规模的综合运用 \important{【专硕·核心】}
\item 平台形态三阶段演进（双边 → 多边 → 生态圈）：划分三维（内部化空间、多元化程度、赋能能力）与转变两动力（协同性提升、动态竞争压力）\important{【专硕·核心】}
\item 平台生态圈的定义、交叉外部性的最大限度内部化与多维交叉锁定 \important{【专硕·核心】}
\item 平台治理三内容：定边界、调关系、平利益（Parker–Van Alstyne 框架）\important{【专硕·核心】}
\item 平台模式与自营模式：开放性、资产轻重与控制权的两维对比 \important{【专硕·扩展】}
\end{enumerate}
\end{chapterbox}

本章揭示了双边市场定价与竞争的完整逻辑：价格结构而非价格水平是平台决策的核心，锁定而非产量是竞争的目标。这一逻辑贯穿了平台从双边形态成长为生态圈的全过程：淘宝以价格结构击败易趣，滴滴与快的以补贴争夺网络规模，阿里与腾讯则把两侧连接扩展为多维度的生态锁定，中国平台经济二十余年的实践正是这一逻辑的连续展开。价格结构的原理至此已经清楚，它在商业层面的具体实现（免费、订阅、佣金、个性化定价）则是下一章的主题。
